\documentclass[10pt]{article}
\usepackage[margin=0.45in]{geometry}
\usepackage[T1]{fontenc}
\usepackage[utf8]{inputenc}
\usepackage{lmodern}
\usepackage{microtype}
\usepackage{amsmath}
\usepackage{booktabs}
\usepackage{enumitem}
\usepackage[hidelinks]{hyperref}
\usepackage{titlesec}

\setlength{\parindent}{0pt}
\setlength{\parskip}{1.5pt}
\setlist[itemize]{leftmargin=1.2em,itemsep=0pt,topsep=1pt}
\setlist[enumerate]{leftmargin=1.4em,itemsep=0pt,topsep=1pt}
\titlespacing*{\section}{0pt}{3pt}{1pt}
\titlespacing*{\subsection}{0pt}{2pt}{1pt}

\titleformat{\section}{\large\bfseries}{\thesection}{0.5em}{}
\titleformat{\subsection}{\normalsize\bfseries}{\thesubsection}{0.5em}{}

\begin{document}
\small

{\Large\bfseries The G\"odel--Kolmogorov Machine: Verifier-Gated Self-Improving Solver Growth for ARC-AGI-3}\hfill {\normalsize Alexander Kolpakov, July 2026}

{\small Code and artifacts: \url{https://github.com/sashakolpakov/gkm} \;·\;
Docs: \url{https://sashakolpakov.github.io/gkm/} \;·\;
Manuscript: \texttt{arc/manuscript/arc\_agi3.tex}}

\textbf{Claim.} The G\"odel--Kolmogorov Machine is a verifier-driven program-growth approach for local ARC-AGI-3 games. A coding proposer grows solver structure, the simulator validates promoted behavior by replay, and the admission loop tries incumbent-leg composition before new code. Retained source states make the resulting acquisition and reuse claims auditable.

The producer is universal across the benchmark: each game uses the same proposer contract, Arena interface, blank scaffold, complexity coordinate, and replay gate. Its learned executable outputs are game- and level-dependent.

\section*{Setting and objective}
The local harness exposes \texttt{step(action) -> frame}, \texttt{levels\_completed}, and \texttt{clone()} for lookahead. Clone-enabled exploration is stronger than the official reset/step interface and is not included in replay action totals. A candidate program is promoted only if fresh replay validates more completed levels under
\[
F(\pi)=R(\pi)+\lambda C(\pi), \qquad R(\pi)=-\operatorname{levels\_completed}(\pi),
\]
where historical $C$ is positive net retained-size growth in the library and player files. Unchanged legs add no charge, but additions and deletions within a file can cancel, so source inspection is required to attribute a low value to reuse.

\section*{What is optimized}
Executable-world-model agents try to build a predictive simulator, verify it, and plan through it. The G\"odel--Kolmogorov Machine treats a world model as only one possible useful structure. The optimized object is broader: solver-program growth. A promoted artifact may be a probe, perception routine, BFS, planner, reusable leg, literal replay path, or world model. The G\"odel--Kolmogorov Machine is therefore closer to a G\"odel/PowerPlay-style self-improving program with a description-length ledger than to a pure model-building agent. We use \textbf{GKM} below only after this full introduction.

\section*{Current promoted artifacts}
\begin{center}
\small
% Generated by arc/build_artifact_docs.py. Do not edit.
\begin{tabular}{@{}lrrr@{}}
\toprule
Game & Levels & Replay actions & Ledger charge \\
\midrule
Frozen v2 (25 games) & 181/183 & 7001 & -- \\
\texttt{wa30} & 9/9 & 597 & 318 \\
\texttt{ls20} & 7/7 & 365 & 760 \\
\bottomrule
\end{tabular}

The frozen release contains one checkpoint record per promoted level. The two displayed game histories additionally contain one uniform manuscript sidecar per level. \texttt{marginal\_C} denotes positive net retained-description growth per source file; additions and deletions within the same file are netted before the positive part, so same-size replacement can receive zero.

\end{center}

\textbf{\texttt{wa30}.} GKM records a logistics game built around carry, helpers, handoffs, neutralisation, and delivery. The solver discovers and preserves reusable structure: freezing target regions before delivered objects overwrite them; complementing autonomous helpers instead of competing with them; exploiting asymmetric carry collision at wall boundaries; neutralising agents that undo delivery; and refactoring repeated transport into ferry legs.

The level-9 WIP trail is the cleanest audit example. A recovered verified suffix was decoded into five repeated grab--carry--release operations and refactored into \texttt{grab\_carry\_release} and \texttt{ferry\_each}. The promoted player is therefore compact composition rather than an opaque replay; the literal suffix is charged, then the repeated structure is compressed into reusable code.

\textbf{\texttt{ls20}.} The uniform acquisition ledger is $40,54,86,114,138,170,158$. The stricter winning-checkpoint audit finds a sharp conditional-AST drop at L2, from 737 to 247 compressed novelty bytes, while the winning player directly calls unchanged \texttt{follow\_cardinal\_runs}. Cumulative executable source does not shrink; the attribution comes from marginal and literal-call evidence together.

\section*{Audit trail}
GKM is meant to leave evidence, not just outputs. Each promoted step is backed by fresh replay, preserved WIP snapshots, a marginal-complexity charge, and a distinction between literal action recovery and reusable refactor. This makes the artifact reviewable: an external evaluator can inspect which information entered through interaction, which code was introduced, what was charged as a literal, and what later became a reusable leg.

\section*{Comparison}
The exact winning-checkpoint test couples conditional AST novelty to direct calls of unchanged definitions. Of 181 admitted boundaries, 180 retain exact historical acquisition source. GKM has 57 direct leg-reuse wins, 14 coupled to half-or-more marginal drops. OPINE has four direct engine-reuse wins, two coupled to sharp drops. baseline1 has no coupled unchanged-world-model witness in its released exact adjacent commands. Retrodict releases curated memory rather than executable winning checkpoints.

\section*{Next evaluation and collaboration request}
The frozen v2 release spans 181/183 levels across all 25 games and 7001 stored actions. Its definitive Competition-Mode replay scores 98.11664037825032\%; the distinct raw coverage is 98.907103825137\%. The next step is a compute-matched comparison against graph exploration and executable world models, plus independent artifact review and private-set evaluation.

\textbf{Bottom line.} GKM treats local solver development as verifier-gated acquisition of reusable program structure. Replay establishes endpoint behavior; checkpoint JSON and retained source expose when novelty was admitted, when prior legs sufficed, and how the solver was refactored.

\end{document}
